\section{Distillation Attempt and Failure Analysis}
\label{sec:distillation}

Phase~1 narrowed the distillation failure to a geometric gap in weight space that SFT does not close. We now test whether this gap reflects a causal bottleneck---and if so, where it lies. We first establish the DPO baseline (\S\ref{sec:dpo_results}), then rule out capacity (\S\ref{sec:capacity}) and evaluation artifacts (\S\ref{sec:truncation}), examine behavioral evidence (\S\ref{sec:suppression}), and test the causal hypothesis (\S\ref{sec:causal}).

%--------------------------------------------------------------
\subsection{Experimental Design}
\label{sec:distill_design}

\paragraph{Training pipeline.}
We train a pooled SFT baseline and three per-$N$ DPO variants on Qwen2.5-3B (LoRA~\cite{hu2022lora}, rank 32). The SFT stage uses 1{,}118 partial-credit trajectories (pooled across $N \in \{2, 4, 8\}$; inclusion criterion: $\text{total\_em} \geq 1$) for 3 epochs with $\text{lr} = 2 \times 10^{-4}$. Each DPO stage initializes from the SFT checkpoint and trains on $N$-specific preference pairs matched at 176 raw pairs ($\approx$109--114 train after hash-based split), with $\beta = 0.05$, $\text{lr} = 2 \times 10^{-5}$, and 1 epoch. Training details and convergence curves are in Appendix~\ref{app:training}.

\paragraph{Matched-sample design.}
The matched-sample constraint is critical for fair cross-$N$ comparison but severely limits data volume. At $N\!=\!4$, the teacher produces only 176 usable preference pairs (from 500 prompts at $15.2\%$ perfect rate); at $N\!=\!8$, only $0.4\%$ of trajectories achieve perfect scores. We subsample all DPO conditions to the $N\!=\!4$ bottleneck of 176 raw pairs. This scarcity is a \emph{structural consequence} of multi-objective collapse: the teacher's own declining success rate at high $N$ limits the supply of contrastive training signal, and generating more prompts does not resolve it.

\paragraph{Evaluation protocol.}
All models are evaluated on held-out test sets (hash-based split, $\text{test\_frac}=0.25$) for $N \in \{2, 4, 8\}$ with $\text{max\_new\_tokens}=4096$, yielding 251, 131, and 125 test prompts respectively. We report em\_partial (primary metric; per-objective EM averaged across objectives and prompts), em\_full (all-$N$-correct), and truncated\_rate (fraction of outputs reaching the token limit). Statistical significance is assessed via paired bootstrap (1{,}000 resamples). We pre-registered a practical-significance threshold of $\Delta \text{em\_partial} \geq +5$pp at $N\!=\!8$ with $p < 0.05$ before Stage~1 training began, representing $\approx$11\% recovery of teacher capability---the minimum we deemed meaningful for justifying the DPO pipeline.

%--------------------------------------------------------------
\subsection{DPO vs.\ SFT: Matched-Sample Comparison}
\label{sec:dpo_results}

Table~\ref{tab:eval_matrix} reports the full evaluation matrix. Under fair comparison (both SFT and DPO evaluated with $\text{max\_new\_tokens}=4096$), DPO provides marginal gains that do not reach practical significance.

\begin{table}[t]
\centering
\small
\resizebox{\columnwidth}{!}{%
\begin{tabular}{lccc}
\toprule
Model & $N\!=\!2$ & $N\!=\!4$ & $N\!=\!8$ \\
\midrule
SFT shared (3B) & 0.112 & 0.084 & 0.025 \\
SFT full-rank (7B)\textsuperscript{$\ddagger$} & 0.171 & 0.147 & 0.079 \\
DPO $N\!=\!2$ & 0.128 & 0.095 & 0.031 \\
DPO $N\!=\!4$ & 0.128 & 0.099 & 0.026 \\
DPO $N\!=\!8$ & 0.130 & 0.082 & 0.035 \\
\midrule
RL teacher (7B) & 0.59 & 0.54 & 0.455 \\
\bottomrule
\end{tabular}%
}
\caption{em\_partial across distillation methods and evaluation conditions ($\text{max\_new\_tokens}=4096$). \textsuperscript{$\ddagger$}The 7B full-rank SFT uses pooled training on $N \in \{2,4,8\}$ (same data and inclusion criteria as the 3B student). The 7B--3B gap accounts for \emph{at most} $\approx$12--14\% of the teacher--student performance gap, remarkably stable across $N$ (\S\ref{sec:capacity}). DPO variants show marginal gains; best $N\!=\!8$: DPO$_{N=8}$ at $+1.0$pp ($p = 0.163$).}
\label{tab:eval_matrix}
\end{table}

\paragraph{Capacity isolation.}
\label{sec:capacity}
To isolate the capacity contribution, we train a 7B full-rank SFT on the same pooled data ($N \in \{2,4,8\}$, identical inclusion criteria) used for the 3B LoRA student. The 7B model consistently outperforms the 3B student: $0.171$ vs.\ $0.112$ at $N\!=\!2$, $0.147$ vs.\ $0.084$ at $N\!=\!4$, and $0.079$ vs.\ $0.025$ at $N\!=\!8$. Defining the capacity contribution as $({\text{em}_{7\text{B}} - \text{em}_{3\text{B}}}) / ({\text{em}_{\text{teacher}} - \text{em}_{3\text{B}}})$, the 7B--3B gap accounts for \emph{at most} $12.3\%$, $13.8\%$, and $12.6\%$ of the teacher--student gap at $N\!=\!2$, $4$, and $8$ respectively---a remarkably stable $\approx$12--14\% upper bound across objective counts. This is an upper bound because the 7B model uses full-rank fine-tuning (more expressive than LoRA), so any capacity advantage is conflated with the training-method advantage. The 7B model also substantially reduces truncation ($16.7\%$ vs.\ $30.7\%$ at $N\!=\!2$; $60.8\%$ vs.\ $83.2\%$ at $N\!=\!8$), indicating that the larger model better controls output length. Nevertheless, the 7B SFT's best performance ($0.171$) remains at $29\%$ of the RL teacher ($0.59$), confirming that capacity is a minor contributing factor while the SFT method itself---specifically, behavioral mismatch (\S\ref{sec:causal})---is the dominant bottleneck. We note that the 7B model uses full-rank fine-tuning while the 3B student uses LoRA (rank 32), so this comparison is a best-effort capacity isolation (\S\ref{sec:limitations}).

The best result---DPO$_{N=8}$ at eval $N\!=\!8$---yields $\Delta = +1.0$pp ($p = 0.163$, 95\% CI $[-0.9, +2.9]$pp), far below the pre-registered practical-significance threshold of $+5$pp. A post-hoc power analysis confirms this is an informative null: with $n = 125$ paired test prompts and observed $\sigma_{\text{diff}} \approx 10.8$pp, we achieve $> 0.99$ power to detect a $5$pp effect and $0.80$ power at $2.7$pp. The non-significant result is not due to insufficient power but reflects a genuinely small effect. We note that the data scarcity is structural rather than remediable (\S\ref{sec:limitations}).

No systematic $N$-matching effect is observed: training on $N\!=\!k$ preference pairs does not specifically benefit evaluation at $N\!=\!k$. DPO$_{N=2}$ performs comparably to DPO$_{N=8}$ across all conditions, suggesting that the preference signal---regardless of objective count---provides similarly limited benefit.

\paragraph{Scale of failure.}
Both students retain $<$8\% of teacher performance at $N\!=\!8$: SFT at $0.025$ ($94.5\%$ drop) and DPO$_{N=8}$ at $0.035$ ($92.3\%$ drop). Neither achieves em\_full $> 0$ at $N \geq 4$.

%--------------------------------------------------------------
\subsection{Truncation Confound}
\label{sec:truncation}

An important methodological episode shaped our analysis. An earlier evaluation used $\text{max\_new\_tokens}=2048$, under which the DPO advantage appeared significant (Table~\ref{tab:truncation_confound}). Re-evaluating SFT with the same 4096-token budget used for DPO revealed that the original significance was largely an artifact.

\begin{table}[t]
\centering
\small
\resizebox{\columnwidth}{!}{%
\begin{tabular}{lcccc}
\toprule
& em\_partial & trunc.\ rate & $\Delta$ vs SFT & $p$ \\
\midrule
\multicolumn{5}{l}{\textit{SFT at max\_new\_tokens = 2048 (confounded):}} \\
\quad SFT shared & 0.011 & 93.6\% & --- & --- \\
\quad DPO $N\!=\!8$ & 0.035 & ---\textsuperscript{$\dagger$} & $+2.4$pp & $<0.001$ \\
\midrule
\multicolumn{5}{l}{\textit{SFT at max\_new\_tokens = 4096 (corrected):}} \\
\quad SFT shared & 0.025 & 83.2\% & --- & --- \\
\quad DPO $N\!=\!8$ & 0.035 & $\sim$84\% & $+1.0$pp & $0.163$ \\
\bottomrule
\end{tabular}%
}
\caption{Truncation confound at $N\!=\!8$. The initial SFT evaluation at 2048 tokens truncated 93.6\% of outputs, yielding a deceptively significant DPO advantage. Extending SFT's budget to 4096 tokens reduced truncation to 83.2\%, lifted em\_partial from 0.011 to 0.025, and revealed that the DPO advantage is not significant. \textsuperscript{$\dagger$}DPO was evaluated at $\text{max\_new\_tokens}=4096$ in both conditions; only the SFT budget differed.}
\label{tab:truncation_confound}
\end{table}

At 2048 tokens, $93.6\%$ of SFT outputs were truncated before the \texttt{<answer>} tag; extending to 4096 allowed $\sim$10\% more to complete, lifting em\_partial from $0.011$ to $0.025$. DPO's score was unchanged at $0.035$---its apparent advantage came from SFT being disproportionately penalized. This confound is specific to RL-distilled models whose procedural strategies produce verbose outputs; we report it as a methodological caution for future work.

%--------------------------------------------------------------
\subsection{Behavioral Mismatch}
\label{sec:suppression}

The truncation confound points to a deeper issue: \emph{why} are the student's outputs so long that $83\%$ hit the 4096-token limit at $N\!=\!8$? The answer reveals a behavioral mismatch between what the student learns and what it can execute.

\paragraph{Observational evidence.}
The evidence in this subsection is \emph{observational}, not causal (\S\ref{sec:limitations}). The SFT student reproduces the teacher's multi-turn action structure (\texttt{<think>} $\to$ \texttt{<search>} $\to$ \texttt{<information>} $\to$ \texttt{<answer>}), but without the retrieval API, search and information segments contain hallucinated content (50--200 tokens per block). At $N\!=\!8$, these hallucinated search cycles consume the majority of the 4096-token budget before the first \texttt{<answer>} tag.

\paragraph{Truncation scaling with $N$.}
Truncation scales monotonically: $30.7\%$ at $N\!=\!2$, $55.7\%$ at $N\!=\!4$, $83.2\%$ at $N\!=\!8$---consistent with generation length growing proportionally to inherited search procedures.

Critically, DPO does not alleviate this pattern. The DPO models exhibit comparable truncation rates ($\sim$84\% at $N\!=\!8$), confirming that preference learning does not teach the student to adopt a more token-efficient strategy---it inherits the same procedural behavior as SFT.\footnote{A pilot attempt to suppress search tokens via bigram blocking was inconclusive: the model generated tag variants (e.g., malformed closing tags, inserted whitespace) that bypassed token-level suppression, suggesting that the behavioral pattern is encoded at a level above individual token sequences.}

%--------------------------------------------------------------
\subsection{Causal Verification: Oracle Retrieval}
\label{sec:causal}

The observational evidence in \S\ref{sec:suppression} raises a natural intervention: if the student's failure stems from executing search procedures without a retrieval backend, then \emph{providing} the backend should recover performance. We test this by intercepting the student's \texttt{<search>} outputs at inference time and injecting oracle retrieval results from the RL teacher's trajectories (the same passages the teacher used for each test prompt).

\begin{table}[t]
\centering
\small
\resizebox{\columnwidth}{!}{%
\begin{tabular}{lccc}
\toprule
& $N\!=\!2$ & $N\!=\!4$ & $N\!=\!8$ \\
\midrule
SFT (no API) & 0.112 & 0.084 & 0.025 \\
SFT + oracle API & \textbf{0.530} & \textbf{0.521} & \textbf{0.498} \\
$\Delta$ & $+41.8$pp & $+43.7$pp & $+47.3$pp \\
$p$ (bootstrap) & ${<}0.0001$ & ${<}0.0001$ & ${<}0.0001$ \\
\midrule
RL teacher\textsuperscript{$\dagger$} & 0.59 & 0.54 & 0.455 \\
\bottomrule
\end{tabular}%
}
\caption{Effect of providing oracle retrieval results at inference time. With real search results injected, the SFT student recovers from near-zero to near-teacher performance. Paired bootstrap, 10{,}000 resamples. \textsuperscript{$\dagger$}Teacher evaluated on a different prompt set; direct comparison is not valid.}
\label{tab:oracle_api}
\end{table}

Table~\ref{tab:oracle_api} shows that providing real retrieval results transforms student performance: em\_partial jumps from $0.025$ to $0.498$ at $N\!=\!8$ ($+47.3$pp, $p < 0.0001$, 95\% CI $[43.8, 50.8]$pp). Truncation drops from $83.2\%$ to $0\%$---the student completes every output when search queries return real passages instead of hallucinated content. This constitutes causal evidence that behavioral mismatch, not capacity limitations alone, drives the distillation failure: the student has learned a viable search strategy but lacks the execution environment to realize it.

\paragraph{Realistic retrieval: BM25.}
The oracle setting uses teacher passages and thus provides an upper bound. To assess a deployable configuration, we replace oracle retrieval with BM25 (top-4 passages from the same corpus). Table~\ref{tab:bm25} and Figure~\ref{fig:retrieval_recovery} show that BM25 retrieval recovers a substantial fraction of the oracle gain: em\_partial improves from $0.025$ to $0.147$ at $N\!=\!8$ ($+12.2$pp), with truncation dropping to $0\%$ across all $N$. The BM25 student recovers $14$--$26\%$ of the oracle improvement, with the largest relative gain at $N\!=\!8$ where the no-API baseline is most degraded. The gap between BM25 and oracle ($0.147$ vs.\ $0.498$ at $N\!=\!8$) reflects retrieval quality rather than student capability: the student's queries closely match the teacher's (Jaccard $= 0.95$), but BM25's top-4 recall is substantially lower than oracle injection.

\begin{table}[t]
\centering
\small
\resizebox{\columnwidth}{!}{%
\begin{tabular}{lccc}
\toprule
& $N\!=\!2$ & $N\!=\!4$ & $N\!=\!8$ \\
\midrule
SFT (no API) & 0.112 & 0.084 & 0.025 \\
SFT + BM25 & 0.203 & 0.147 & 0.147 \\
SFT + oracle API & \textbf{0.530} & \textbf{0.521} & \textbf{0.498} \\
\midrule
RL teacher & 0.59 & 0.54 & 0.455 \\
\bottomrule
\end{tabular}%
}
\caption{em\_partial under different retrieval conditions. BM25 retrieval recovers a meaningful fraction of the oracle gain ($+9.2$pp at $N\!=\!2$, $+12.2$pp at $N\!=\!8$) and eliminates truncation entirely ($0\%$ across all $N$). The remaining gap to oracle reflects retrieval quality, not student strategy.}
\label{tab:bm25}
\end{table}

\begin{figure}[t]
\centering
\includegraphics[width=\columnwidth]{retrieval_recovery.pdf}
\caption{em\_partial under different retrieval conditions. Providing oracle retrieval results at inference time recovers the SFT student from near-zero to near-teacher performance across all $N$, while even a standard BM25 backend yields meaningful gains. The RL teacher (hatched) is shown for reference.}
\label{fig:retrieval_recovery}
\end{figure}

\paragraph{Query fidelity and search behavior.}
The student averages exactly $1.0$ search per prompt across all 507 test cases. Crucially, this is not a simplification of the teacher's strategy: analysis of the teacher's training trajectories reveals that $>$95\% of teacher responses also contain a single search round (mean searches: $1.04$ at $N\!=\!2$, $1.01$ at $N\!=\!4$, $1.06$ at $N\!=\!8$), because the prompt prefix already includes the first retrieval result. Furthermore, the student's search queries closely match the teacher's: mean Jaccard similarity is $0.95$ (median $= 1.0$, $n = 507$). SFT thus faithfully transfers both the search pattern and the query formulation---the bottleneck is purely the absence of a retrieval backend, not a failure of strategy learning.

\paragraph{Why em\_full remains low.}
Despite recovering em\_partial to near-teacher levels, em\_full remains low under oracle retrieval ($0.008$ at $N\!=\!8$). Per-objective analysis reveals that individual objective accuracy is $\approx$50\% (range 38--58\% across positions), with no systematic decay across objective index. The low em\_full follows from probability stacking: if each of $N$ objectives is independently correct with probability $\approx 0.5$, the probability of \emph{all} $N$ being correct is $\approx 0.5^N$ ($0.4\%$ at $N\!=\!8$), consistent with the observed $0.8\%$. The single search round retrieves information relevant to a subset of objectives; full coverage would require multiple targeted searches.

%--------------------------------------------------------------
\subsection{Synthesis}
\label{sec:synthesis}

The analyses in this section converge on a consistent diagnostic:

\begin{enumerate}[leftmargin=*,nosep]
    \item \textbf{Preference learning does not close the geometric gap.} RL weight changes are $2$--$6\times$ supra-random aligned with the base model's principal directions (\S\ref{sec:svd}), while SFT achieves only $1.47\times$. DPO---which provides richer training signal---does not bridge this gap: all DPO variants remain within $1$pp of SFT at $N\!=\!8$.

    \item \textbf{The student imitates a strategy it cannot execute---causally supported.} The RL teacher's multi-objective capability is inseparable from its iterative search procedure. SFT transfers this procedure, but without the retrieval backend, hallucinated search traces exhaust the generation budget. Providing oracle retrieval at inference time recovers em\_partial from $0.025$ to $0.498$ at $N\!=\!8$ ($+47.3$pp, $p < 0.0001$; \S\ref{sec:causal}), confirming that the student possesses a viable search strategy but lacks execution capability.

    \item \textbf{Truncation masks and distorts the comparison.} At insufficient generation budgets, truncation disproportionately penalizes the more verbose model (SFT), creating a spurious advantage for DPO. After controlling for this confound, the DPO advantage shrinks from $+2.4$pp ($p < 0.001$) to $+1.0$pp ($p = 0.163$)---both methods fail for the same reason.
\end{enumerate}

The oracle retrieval experiment (\S\ref{sec:causal}) elevates behavioral mismatch from observational hypothesis to a causally supported finding, and the BM25 experiment confirms that even a standard retrieval backend recovers meaningful performance ($+12.2$pp at $N\!=\!8$). Notably, the student faithfully reproduces the teacher's search strategy (Jaccard $= 0.95$; \S\ref{sec:causal}), confirming that the bottleneck is execution capability rather than strategy learning. A 7B pooled SFT (\S\ref{sec:capacity}) quantifies the capacity contribution at \emph{at most} $\approx$12--14\% of the teacher--student gap---a minor but consistent factor across all $N$ values, while behavioral mismatch accounts for the remaining $\approx$86--88\% (Figure~\ref{fig:gap_decomposition}).

\begin{figure}[t]
\centering
\includegraphics[width=\columnwidth]{gap_decomposition.pdf}
\caption{Decomposition of the teacher--student performance gap across objective counts. Capacity (7B vs.\ 3B SFT) accounts for a stable ${\sim}$12--14\%, while behavioral mismatch dominates at ${\sim}$86--88\%.}
\label{fig:gap_decomposition}
\end{figure} We note that capacity and behavioral mismatch may interact non-additively, and additional factors (e.g., training-method differences) may also contribute; these percentages should be interpreted as approximate decompositions rather than exact attributions.
